% ============================================================
%  EXAMEN TEMA 2: RAZONES TRIGONOMÉTRICAS
%  Curso 4to Año — Matemática
%  Duración: 90 minutos
%  Temas: Razones en el triángulo rectángulo, ángulos notables
%         y resolución de triángulos rectángulos
%  Autor: Ing. Gabriel Astudillo
%  Nota: enunciado limpio (sin pistas ni desarrollo). Las
%  soluciones están en la "Clave de Respuestas" al final.
% ============================================================

\documentclass[12pt, letterpaper]{article}

% ── Paquetes ─────────────────────────────────────────────────

\usepackage{mathwizards-guia}
\mwguia{Examen — Tema 2: Razones Trigonométricas}{Ing. Gabriel Astudillo $\cdot$ 4to Año}

% ── Comandos propios de este examen ───────────────────
\providecommand{\sen}{\sin}

\begin{document}
% ============================================================

% ── Portada / Título ─────────────────────────────────────────
\mwportada{Examen del Tema 2}{Razones Trigonométricas}{Seno, Coseno y Tangente: Triángulos Rectángulos y Aplicaciones}

\vspace{4pt}
\begin{cuadroinfo}
  \small
  \textbf{Nombre:} \rule{0.55\linewidth}{0.4pt} \quad \textbf{Fecha:} \rule{0.15\linewidth}{0.4pt} \\[6pt]
  \textbf{Tiempo disponible:} 90 minutos \quad$\bullet$\quad \textbf{Puntaje total:} 100 puntos \\[4pt]
  \textbf{Instrucciones:}
  \begin{enumerate}[leftmargin=*]
    \item Lee cada ejercicio con calma antes de resolverlo.
    \item Dibuja un esquema en los problemas de aplicación e identifica el triángulo rectángulo.
    \item Usa los valores exactos de los ángulos notables (no aproximaciones).
    \item Racionaliza los denominadores cuando sea necesario.
    \item No se permite el uso de calculadora.
    \item Escribe con letra clara y revisa tu examen antes de entregarlo.
  \end{enumerate}
\end{cuadroinfo}

\vspace{8pt}

% ============================================================
\section{Razones Trigonométricas en el Triángulo Rectángulo (32 puntos)}
% ============================================================

\begin{enumerate}[leftmargin=*]
  \item \textbf{(8 pts)} Un triángulo rectángulo tiene hipotenusa $25$ y un cateto de $7$.
        \begin{enumerate}[label=\alph*)]
          \item Halla el otro cateto (aplica Pitágoras).
          \item Escribe las \textbf{seis} razones trigonométricas del ángulo agudo $\theta$ opuesto al cateto de $7$.
        \end{enumerate}

  \item \textbf{(8 pts)} Si $\sen \alpha = \dfrac{5}{13}$ y $\alpha$ es un ángulo agudo, calcula $\cos \alpha$ y $\tan \alpha$ usando la identidad fundamental.

  \item \textbf{(8 pts)} Verifica la identidad fundamental $\sen^2 \theta + \cos^2 \theta = 1$ para el ángulo agudo del triángulo de lados $8$, $15$ y $17$.

  \item \textbf{(8 pts)} El cateto opuesto a un ángulo agudo $\alpha$ mide $6$ y la hipotenusa mide $10$. Halla $\sen \alpha$, $\cos \alpha$ y $\tan \alpha$.
\end{enumerate}

\newpage

% ============================================================
\section{Ángulos Notables: $30^{\circ}$, $45^{\circ}$ y $60^{\circ}$ (34 puntos)}
% ============================================================

\begin{enumerate}[leftmargin=*]
  \item \textbf{(9 pts)} Calcula con los valores exactos de la tabla:
        \[
          a)\ \sen 30^{\circ} \qquad b)\ \cos 45^{\circ} \qquad c)\ \tan 60^{\circ} \qquad d)\ \sen 60^{\circ} \qquad e)\ \cos 30^{\circ}
        \]

  \item \textbf{(8 pts)} En un triángulo rectángulo la hipotenusa mide $16$ y uno de los ángulos agudos mide $30^{\circ}$. Halla la medida de los dos catetos.

  \item \textbf{(8 pts)} Un triángulo rectángulo isósceles tiene catetos de $5$\,cm. \textquestiondown Cuánto mide su hipotenusa? \textquestiondown Cuánto miden sus ángulos agudos?

  \item \textbf{(9 pts)} Calcula y explica el resultado:
        \begin{enumerate}[label=\alph*)]
          \item $\tan 30^{\circ} \cdot \tan 60^{\circ}$
          \item $\sen^2 30^{\circ} + \cos^2 30^{\circ}$
        \end{enumerate}
\end{enumerate}

\newpage

% ============================================================
\section{Resolución de Triángulos Rectángulos y Aplicaciones (34 puntos)}
% ============================================================

\begin{enumerate}[leftmargin=*]
  \item \textbf{(9 pts)} Resuelve el triángulo rectángulo cuyos catetos miden $9$\,cm y $12$\,cm (halla la hipotenusa y los dos ángulos agudos).

  \item \textbf{(8 pts)} Desde un punto a $40$\,m de la base de un edificio, el ángulo de elevación hacia la azotea es de $60^{\circ}$. \textquestiondown Qué altura tiene el edificio?

  \item \textbf{(9 pts)} Una escalera de $8$\,m apoyada contra una pared forma un ángulo de $45^{\circ}$ con el suelo.
        \begin{enumerate}[label=\alph*)]
          \item \textquestiondown A qué altura de la pared llega la escalera?
          \item \textquestiondown A qué distancia de la pared está la base de la escalera?
        \end{enumerate}

  \item \textbf{(8 pts)} Un poste de $5$\,m proyecta una sombra de $5\sqrt{3}$\,m. \textquestiondown Cuál es el ángulo de elevación del sol?
\end{enumerate}

\newpage

% ============================================================
\ifmwclaves
  \section{Clave de Respuestas}
  % ============================================================

  \subsection*{Sección 1: Razones Trigonométricas en el Triángulo Rectángulo}

  \begin{enumerate}[leftmargin=*, label=\textbf{\arabic*.}]
    \item a) $b = \sqrt{25^2 - 7^2} = \sqrt{625 - 49} = \sqrt{576} = 24$.
          \quad b) Con opuesto $= 7$, adyacente $= 24$ e hipotenusa $= 25$:
          \[
            \sen \theta = \frac{7}{25}, \quad \cos \theta = \frac{24}{25}, \quad \tan \theta = \frac{7}{24}, \quad \csc \theta = \frac{25}{7}, \quad \sec \theta = \frac{25}{24}, \quad \cot \theta = \frac{24}{7}.
          \]

    \item $\cos^2 \alpha = 1 - \left(\dfrac{5}{13}\right)^2 = 1 - \dfrac{25}{169} = \dfrac{144}{169}$. Como $\alpha$ es agudo, $\cos \alpha = \dfrac{12}{13}$.
          Entonces $\tan \alpha = \dfrac{\sen \alpha}{\cos \alpha} = \dfrac{5/13}{12/13} = \dfrac{5}{12}$.

    \item Tomamos el ángulo $\theta$ con opuesto $8$ e hipotenusa $17$: $\sen \theta = \dfrac{8}{17}$ y $\cos \theta = \dfrac{15}{17}$.
          \[
            \sen^2 \theta + \cos^2 \theta = \left(\frac{8}{17}\right)^2 + \left(\frac{15}{17}\right)^2 = \frac{64 + 225}{289} = \frac{289}{289} = 1. \ \checkmark
          \]

    \item El cateto adyacente mide $\sqrt{10^2 - 6^2} = 8$. Entonces:
          \[
            \sen \alpha = \frac{6}{10} = \frac{3}{5}, \qquad \cos \alpha = \frac{8}{10} = \frac{4}{5}, \qquad \tan \alpha = \frac{6}{8} = \frac{3}{4}.
          \]
  \end{enumerate}

  \newpage

  \subsection*{Sección 2: Ángulos Notables}

  \begin{enumerate}[leftmargin=*, label=\textbf{\arabic*.}]
    \item a) $\sen 30^{\circ} = \dfrac{1}{2}$ \quad b) $\cos 45^{\circ} = \dfrac{\sqrt{2}}{2}$ \quad c) $\tan 60^{\circ} = \sqrt{3}$ \quad d) $\sen 60^{\circ} = \dfrac{\sqrt{3}}{2}$ \quad e) $\cos 30^{\circ} = \dfrac{\sqrt{3}}{2}$

    \item Cateto opuesto al ángulo de $30^{\circ}$: $\sen 30^{\circ} = \dfrac{x}{16} \Rightarrow x = 16 \cdot \dfrac{1}{2} = 8$.
          Cateto adyacente: $\cos 30^{\circ} = \dfrac{y}{16} \Rightarrow y = 16 \cdot \dfrac{\sqrt{3}}{2} = 8\sqrt{3}$.

    \item Hipotenusa: $c = \sqrt{5^2 + 5^2} = \sqrt{50} = 5\sqrt{2}$\,cm.
          Los ángulos agudos son iguales (triángulo isósceles) y suman $90^{\circ}$: cada uno mide $45^{\circ}$.

    \item a) $\tan 30^{\circ} \cdot \tan 60^{\circ} = \dfrac{\sqrt{3}}{3} \cdot \sqrt{3} = \dfrac{3}{3} = 1$. (Los ángulos son complementarios.)
          \quad b) $\sen^2 30^{\circ} + \cos^2 30^{\circ} = \left(\dfrac{1}{2}\right)^2 + \left(\dfrac{\sqrt{3}}{2}\right)^2 = \dfrac{1}{4} + \dfrac{3}{4} = 1$. (Es la identidad fundamental.)
  \end{enumerate}

  \newpage

  \subsection*{Sección 3: Resolución de Triángulos Rectángulos y Aplicaciones}

  \begin{enumerate}[leftmargin=*, label=\textbf{\arabic*.}]
    \item Hipotenusa: $c = \sqrt{9^2 + 12^2} = \sqrt{81 + 144} = \sqrt{225} = 15$\,cm.
          Ángulo $\alpha$ frente al cateto de $9$: $\tan \alpha = \dfrac{9}{12} = 0{,}75 \Rightarrow \alpha \approx 36{,}9^{\circ}$.
          El otro ángulo: $\beta = 90^{\circ} - 36{,}9^{\circ} \approx 53{,}1^{\circ}$.

    \item $\tan 60^{\circ} = \dfrac{h}{40} \Rightarrow \sqrt{3} = \dfrac{h}{40} \Rightarrow h = 40\sqrt{3} \approx 69{,}3$\,m.

    \item a) Altura: $\sen 45^{\circ} = \dfrac{h}{8} \Rightarrow \dfrac{\sqrt{2}}{2} = \dfrac{h}{8} \Rightarrow h = 4\sqrt{2} \approx 5{,}66$\,m.
          \quad b) Distancia de la base: como el ángulo es $45^{\circ}$, ambos catetos son iguales, así que la distancia también es $4\sqrt{2} \approx 5{,}66$\,m
          (o bien $\cos 45^{\circ} = \dfrac{d}{8} \Rightarrow d = 4\sqrt{2}$).

    \item $\tan \theta = \dfrac{5}{5\sqrt{3}} = \dfrac{1}{\sqrt{3}} = \dfrac{\sqrt{3}}{3} = \tan 30^{\circ}$. Por lo tanto, el ángulo de elevación del sol es $\theta = 30^{\circ}$.
  \end{enumerate}

\fi
\end{document}
